\begin{table}[ht]
\centering
\caption{Summary Statistics}
\label{tab:summary_stats}
\begin{threeparttable}
\begin{tabular}{l>{\hspace{1.5em}}r>{\hspace{1.5em}}r>{\hspace{1.5em}}r>{\hspace{1.5em}}r}
\hline\hline
 & Obs. & Mean & Std. Dev. & Max \\
\hline
\multicolumn{5}{l}{\textit{Trade value (million USD)}} \\
\quad All         & 6,912 & 6,295.88 & 40,783.45 & 1,279,718.11 \\
\quad Intrastate    & 144 & 164,661.73 & 229,890.30 & 1,279,718.11 \\
\quad Interstate & 6,768 & 2,926.40 & 6,068.23 & 73,050.55 \\
\hline
\multicolumn{5}{l}{\textit{Distance (miles)}} \\
\quad All         & 6,912 & 1,191.24 & 704.91 & 3,130.00 \\
\quad Intrastate    & 144 & 138.29 & 104.84 & 433.59 \\
\quad Interstate & 6,768 & 1,213.64 & 695.08 & 3,130.00 \\
\hline\hline
\end{tabular}

\begin{tablenotes}[online,flushleft]
    \small
    \item \textit{Notes:} Data on bilateral trade values for the lower 48 U.S. states are for the years 2012, 2017, and 2022, and are taken from the FAF5 (2012 and 2017) and FAF6 (2022) datasets from the Bureau of Transportation Statistics. Data on bilateral distances are calculated according to the procedure described in the text. Distances refer to least-cost diesel routes; electric route distances, defined only for charging-feasible pairs, are summarized in Table~\ref{tab:network_development}.
 \end{tablenotes}
 
\end{threeparttable}
 \end{table}